\documentclass[11pt]{article}
\usepackage[a4paper,margin=25mm]{geometry}
\usepackage{amsmath,amssymb,hyperref}
\hypersetup{colorlinks=true,urlcolor=blue}
\title{Conjecture 469 is false:\\a triangle and its plane dual share the root 1}
\author{GodBlf}
\date{October 4, 2026}
\begin{document}
\maketitle

\section{Claim and counterexample}
Conjecture 00000000469 asserts that the flow and tension polynomials of
a pair of planar dual graphs have disjoint sets of real roots, with a
suggested proof using sign reversal on dual cycles. Under the standard
nowhere-zero definitions this is false.

Let $G=C_3$ be a triangle embedded in the plane. It has two faces, the
bounded face and the exterior face. Each of its three edges separates
these faces. Therefore its plane dual $H=G^*$ has two vertices joined
by three parallel edges. Parallel edges are part of the standard plane
dual construction and are not excluded by the conjecture.
We will prove
\[
             F_G(q)=T_H(q)=q-1.
\]
Thus their real-root sets are both $\{1\}$, refuting the claimed separation.

\section{Definitions and direct calculation}
For a finite graph $X$, choose arbitrary edge orientations and a finite
abelian group $\Gamma$ of order $q$. A nowhere-zero flow assigns a nonzero
group element to each edge so that incoming and outgoing sums agree at
every vertex. A nowhere-zero tension is an edge assignment of the form
$p(v)-p(u)$ on the oriented edge $u\to v$, with all edge values nonzero.
The number of these assignments gives $F_X(q)$ and $T_X(q)$, respectively.
Reversing an edge orientation negates its value and preserves these counts.

Orient the three edges of $G$ cyclically. The conservation equations force
all three values to be the same element $a\in\Gamma$. Every nonzero $a$
works, so $F_G(q)=q-1$ for every positive integer $q$.
Orient all three edges of $H$ from its first vertex to its second.
All tension values equal the same potential difference $b$. Each nonzero
$b$ is attainable, and potentials differing by a constant give the same
assignment. Hence $T_H(q)=q-1$. Equality at infinitely many positive
integers identifies the polynomials, and their evaluation at the real
number $1$ is zero. In particular, $q=1$ is legitimate: the trivial group
has no nonzero element, so there are no nowhere-zero assignments.

For completeness, the other two polynomials are
\[
 T_G(q)=(q-1)(q-2),\qquad F_H(q)=(q-1)(q-2).
\]
Indeed $T_G=P_G/q$, with $P_G=q(q-1)(q-2)$. A flow on $H$ is a triple
of nonzero elements $(a,b,c)$ satisfying $a+b+c=0$: choose $a\ne0$,
then $b\ne0,-a$, and $c$ is forced. Thus either ordering of the dual pair,
and even comparison on the same triangle, gives a common root at $1$.

\section{Why planar duality contradicts separation}
The standard planar duality identities are
\[
 F_X(q)=T_{X^*}(q),\qquad T_X(q)=F_{X^*}(q).
\]
Equivalently, for a connected plane graph,
$F_X(q)=P_{X^*}(q)/q$ and $T_X(q)=P_X(q)/q$.
These are polynomial identities; the quotients are polynomials, with
no evaluation by division at $q=0$. The cycle/cut correspondence converts
flow constraints into tension constraints, not separation of roots.
Identical polynomials have identical root sets, so any real root suffices.
If one additionally imposed simple graphs on both sides, the usual
self-dual tetrahedral embedding of $K_4$ would also refute the assertion:
both polynomials equal $(q-1)(q-2)(q-3)$. No such extra restriction is
needed for the triangle example.

\section{Formal verification and its scope}
Write $k_X(A)$ for the number of connected components of the spanning
subgraph with edge set $A\subseteq E$, including isolated vertices.
The standard subset expansions are
\begin{align*}
F_X(q)&=\sum_{A\subseteq E}(-1)^{|E|-|A|}
                    q^{|A|-|V|+k_X(A)},\\
T_X(q)&=\sum_{A\subseteq E}(-1)^{|A|}
                    q^{k_X(A)-k_X(E)}.
\end{align*}
They follow from inclusion--exclusion for nonzero edge values.
All exponents are nonnegative. The tension expansion is also
$P_X(q)/q^{k_X(E)}$. At $q=1$, every power is one, and the two sums
for three edges become respectively $-1+3-3+1=0$ and $1-3+3-1=0$.

The Lean project defines graphs by their actual endpoint maps. Its
\texttt{components} counts equivalence classes of the equivalence closure
of the selected edge adjacency relation, rather than inserting a table
of assumed component counts. The functions \texttt{flow} and
\texttt{tension} use the expansions above over $\mathbb R$. Theorems
\texttt{flow\_at\_one} and \texttt{tension\_at\_one} prove the zero
evaluation for every graph with three labeled edges, so no computation
of quotient cardinalities is required. \texttt{common\_real\_root}
and \texttt{conjecture\_false} specialize this to the actual triangle
and triple bond.

The six darts of the triangle also carry an explicit rotation system.
Lean checks vertex preservation, the rotation involution, the two
three-dart face cycles, their exhaustion, connectivity of both graphs,
the Euler count $3-3+2=2$, and the dual endpoint maps. The interpretation
of this rotation system as a plane embedding, and the standard
inclusion--exclusion correspondence to group counts, are justified in
this report; Lean does not formalize the surface classification theorem
or the group-counting derivation. Its root theorems use the standard
subset definitions directly, not postulated polynomial formulas.

The auxiliary Python script independently computes component counts
and coefficient arrays from the same subset definitions, and compares
them with the four explicit polynomials above. Neither this computation
nor a floating-point root search is needed for the mathematical disproof.
The Lean project uses no admitted proofs, additional axioms, or
\texttt{native\_decide}; the axiom audit is included in
\texttt{verification.txt}.

\end{document}
